% generated by roam.py handbook export -- не редактировать, править заметки
% 20261014120524-косинус_угла_между_градиентом_и_направлением_спуска.org (body bd709cd58c67)

\begin{definition}[{Косинус угла между градиентом и направлением спуска}]\label[definition]{note:E4AE9B25-BDB9-4A54-8CCB-09809F20B79D}
  Вычисляется по формуле косинуса между \hyperref[note:1756907B-1FF2-4F62-8553-759B8509281C]{градиентом} $\nabla f(x_k) \ne 0$ и \hyperref[note:BEEE81AD-13FC-4846-A160-C72E31760BEE]{направлением спуска} $d_k \ne 0$
  \[
  \cos \theta_k = \frac{\left\langle \nabla f(x_k); d_k \right\rangle}{\|\nabla f(x_k)\| \cdot \|d_k\|}
  \]
\end{definition}

\begin{remark}
  Для направления спуска $\cos \theta_k < 0$, для антиградиента $\cos \theta_k = -1$. Условие $|\cos \theta_k| > C$ в \hyperref[note:36F6E1BA-41DB-460D-A9C6-B980621727D4]{сходимости градиентного спуска с шагом Армихо-Вульфа} запрещает направлениям становиться почти ортогональными градиенту; для метода Ньютона оно проверяется в \hyperref[note:CC41F7D2-EECB-46E4-817A-756E7CB9B7C0]{глобальной сходимости}: $\cos^2 \theta_k \ge \frac{\mu^2}{L^2}$.
\end{remark}
